\begin{table}[htbp]
\centering\small
\caption{分块到达顺序下的期望收入}\label{tab:order}
\begin{tabular}{lrrrrrrr}
\toprule
到达顺序 & \Greedy{} & \LPFD{} & \LPDual{} & \LPCount{} & \IIDDP{} & \SeqDP{} & $\E c^{\LPDual}_{T+1,1}$ \\
\midrule
$1\to2\to3$ & 694.050 & 694.200 & 694.199 & 690.984 & 693.735 & 694.543 & 0.851 \\
$1\to3\to2$ & 693.395 & 693.626 & 693.588 & 690.307 & 693.126 & 693.946 & 0.923 \\
$2\to1\to3$ & 696.997 & 674.747 & 683.029 & 696.999 & 670.601 & 697.000 & 2.360 \\
$2\to3\to1$ & 696.806 & 694.261 & 695.511 & 696.782 & 694.757 & 696.812 & 0.830 \\
$3\to1\to2$ & 693.395 & 693.444 & 693.416 & 691.694 & 693.588 & 693.908 & 0.923 \\
$3\to2\to1$ & 696.806 & 696.075 & 696.707 & 696.807 & 696.498 & 696.809 & 0.923 \\
\bottomrule
\end{tabular}
\par\vspace{0.35em}{\footnotesize 表中为精确期望收入，各顺序的 $F_s$ 均为 706。$1\to2\to3$ 表示先到达 50 位类型 1 顾客，再到达 30 位类型 2 与 20 位类型 3 顾客，其余类推。算法 1 在任何顺序下的期望收入都是 \AlgOneER{}。}
\end{table}
